\documentclass{article}

\usepackage{parskip}

\author{Stefan Ehin}
\title{Notes on: Dynamical Low Rank approximation of PDEs with random parameters}

\begin{document}

\maketitle

\section*{Problem}
\begin{enumerate}
    \item Regular PDE problem statements do not account for the potential error in many parameters: coefficients of the operator, forcing/source terms, initial conditions and boundary conditions.
    \item
    \begin{enumerate}
        \item Monte Carlo method is simple but has slow convergence in sample size.
        \item Spectral methods are considerably more compute-efficient but it assumes the underlying solution map is smooth and the computation becomes compute-heavy when the stochastic space is high-dimensinal.
        \item In case the solution manifold is approximately linear the curse of dimensionality is broken but it fails for manifolds that strongly dimiate from being linear.
    \end{enumerate}
\end{enumerate}

\section*{Key idea}
\begin{enumerate}
    \item Model all the error-prone variables as probability distributions.
    \item Use a dynamical low-rank basis i.e. reduced basis like in 2. (C) but evolving over time.
\end{enumerate}

\section*{What I'd prove/extend}
	
\end{document}

